\section{Constitutive models implemented in the research solver}
\label{sec:constitutive}
The paste law combines isotropic Hencky elasticity and a principal-frame
viscoplastic return, following the general class of Herschel-Bulkley MPM models
used for yield-stress materials \cite{foam}. The liquid law is a weakly
compressible Newtonian approximation. The equations below describe this code,
not all features of the cited methods and not measured fecal material.

\subsection{Elastic constants, stresses, and shear-rheometer conventions}
From Young's modulus $E$ and Poisson ratio $\nu_P$,
\begin{equation}
 G=\frac{E}{2(1+\nu_P)},\qquad
 \lambda_L=\frac{E\nu_P}{(1+\nu_P)(1-2\nu_P)},\qquad
 K_b=\lambda_L+\frac{2G}{3}=\frac{E}{3(1-2\nu_P)}.
 \label{eq:moduli}
\end{equation}
The symbol $\lambda_L$ denotes the first Lam\'e parameter, not a principal
stretch or interpolation weight. The wire contract uses rheometer shear values,
\begin{equation}
 \tau_{\rm shear}=\tau_y+K_s|\dot\gamma|^n,
\end{equation}
whereas the internal J2 law uses an equivalent deviatoric stress $q$ and
equivalent plastic strain rate $\dot\varepsilon_p$. In simple shear,
\begin{equation}
 q=\sqrt{3}\,|\tau_{\rm shear}|,\qquad
 \dot\varepsilon_p=\frac{|\dot\gamma_p|}{\sqrt{3}},\qquad
 \boxed{\sigma_y=\sqrt{3}\,\tau_y,\quad K_e=3^{(1+n)/2}K_s}.
 \label{eq:shear-conversion}
\end{equation}
Thus the internal law is $q=\sigma_y+K_e\dot\varepsilon_p^n$ during plastic
flow. Treating the same numeric $K$ as both $K_s$ and $K_e$ would change the
material. Moreover, the code's equivalent stress is formed from
\emph{Kirchhoff} stress. Ordinary rheometer shear stress is a Cauchy stress;
the identification above presumes $J\simeq1$. At finite volume change, Cauchy
equivalent stress is $q/J$, so this is not an exact constant-Cauchy-yield law
under arbitrary compression.

\subsection{Trial Hencky strain and the scalar implicit return}
Let $\mathbf F_e^n$ be the stored elastic tensor and $\mathbf C$ the new affine
velocity gradient from the grid. Form
\begin{align}
 \mathbf F_e^{\rm tr}&=(\mathbf I+\Delta t\mathbf C)\mathbf F_e^n,\\
 \mathbf b^{\rm tr}&=\mathbf F_e^{\rm tr}(\mathbf F_e^{\rm tr})^T
   =\mathbf U\operatorname{diag}(\ell_1^2,\ell_2^2,\ell_3^2)\mathbf U^T,\\
 e_i^{\rm tr}&=\ln\ell_i,\qquad
 \bar e=\tfrac13\sum_i e_i^{\rm tr},\qquad
 \mathbf e_d^{\rm tr}=\mathbf e^{\rm tr}-\bar e\mathbf 1.
 \label{eq:hencky-trial}
\end{align}
The code diagonalizes $\mathbf b^{\rm tr}$ with a cyclic Jacobi eigensolver;
it does not require a separately computed right singular basis. With
$\mathbf s=\operatorname{dev}\boldsymbol\tau$,
\begin{equation}
 q_{\rm tr}=\sqrt{\tfrac32\mathbf s_{\rm tr}:\mathbf s_{\rm tr}}
           =\sqrt{6}\,G\|\mathbf e_d^{\rm tr}\|.
\end{equation}
If $q_{\rm tr}\leq\sigma_y$, the state is elastic. Otherwise the discrete J2
stress reduction and overstress law give
\begin{equation}
 q=q_{\rm tr}-3G\Delta\varepsilon_p,\qquad
 \Delta\varepsilon_p=\Delta t
       \left(\frac{q-\sigma_y}{K_e}\right)^{1/n}.
\end{equation}
Eliminating the increment yields the actual scalar residual
\begin{equation}
 \boxed{f(q)=q-q_{\rm tr}+3G\Delta t
       \left(\frac{q-\sigma_y}{K_e}\right)^{1/n}=0},
 \quad q\in[\sigma_y,q_{\rm tr}]. \label{eq:hb-return}
\end{equation}
For $K_e>0$, $n>0$ and $q_{\rm tr}>\sigma_y$,
$f(\sigma_y)<0$, $f(q_{\rm tr})>0$, and in the bracket interior
\begin{equation}
 f'(q)=1+\frac{3G\Delta t}{nK_e}
       \left(\frac{q-\sigma_y}{K_e}\right)^{1/n-1}>0.
\end{equation}
There is a unique root. The code takes a Newton step only if it is finite and
strictly inside the updated bracket; otherwise it bisects. The stopping test
uses either residual or bracket width below
$10^{-12}\max(q_{\rm tr},\SI{1}{Pa})$, with a 200-iteration cap.
The monotonicity argument explains robustness, but a finite floating-point
iteration can still fail and trigger step rejection. $K_e=0$ instead uses the
perfectly plastic result $q=\sigma_y$ directly.

\subsection{Reconstructing the elastic tensor and estimating plastic loss}
Set $\alpha=q/q_{\rm tr}$ in the yielded branch and $\alpha=1$ in the elastic
branch. Then
\begin{align}
 \mathbf e^{n+1}&=\bar e\mathbf1+\alpha\mathbf e_d^{\rm tr},\\
 \mathbf F_e^{n+1}&=\mathbf U\operatorname{diag}
       \left(\frac{e^{e_1^{n+1}}}{\ell_1},
             \frac{e^{e_2^{n+1}}}{\ell_2},
             \frac{e^{e_3^{n+1}}}{\ell_3}\right)\mathbf U^T\mathbf F_e^{\rm tr},
 \label{eq:hencky-reconstruction}\\
 \boldsymbol\tau^{n+1}&=\mathbf U\operatorname{diag}
       \left(2G e_i^{n+1}+\lambda_L\sum_j e_j^{\rm tr}\right)\mathbf U^T.
\end{align}
The sum of log strains is unchanged by the return, so plastic correction is
isochoric. The stored tensor is the returned \emph{elastic} tensor, not a full
accumulated deformation tensor or a separately stored plastic gradient.
Inverted/nonfinite trial tensors, stretches below $10^{-6}$, eigensolver failure
or nonfinite stress reject the trial rather than substituting an undeformed
particle. The principal-frame return and the explicit kinematic trial remain
discrete constitutive approximations.

The reference-volume Hencky elastic energy used by the diagnostic is
\begin{equation}
 \Psi_e=G\sum_i e_i^2+\frac{\lambda_L}{2}\left(\sum_i e_i\right)^2
       =G\|\mathbf e_d\|^2+\frac{K_b}{2}(\operatorname{tr}\mathbf e)^2.
\end{equation}
The code accumulates the backward-Euler plastic-loss estimate
\begin{equation}
 \Delta D_P=V_0 q\Delta\varepsilon_p\geq0. \label{eq:plastic-loss}
\end{equation}
It is useful to distinguish this from the elastic energy released by the
\emph{isolated radial-return operation}. Since $\Psi_{e,d}=q^2/(6G)$ and the
volumetric part does not change,
\begin{align}
 \Delta\Psi_{e,\rm return}
 &=\frac{q_{\rm tr}^2-q^2}{6G}
   =\tfrac12(q_{\rm tr}+q)\Delta\varepsilon_p,\\
 \Delta\Psi_{e,\rm return}-q\Delta\varepsilon_p
 &=\tfrac32G(\Delta\varepsilon_p)^2. \label{eq:return-energy-gap}
\end{align}
This algebraic difference is nonnegative and quadratic in the increment. It is
not a bound on the unknown full physical work history during an MPM step, nor
does it account for transfers, walls or liquid viscosity. It illustrates why
positive plastic loss alone does not close the global discrete energy budget.

\subsection{Worked Herschel-Bulkley return}
Use the synthetic slump parameters $E=\SI{1000}{Pa}$, $\nu_P=0.2$,
$\tau_y=\SI{30}{Pa}$, $K_s=\SI{10}{Pa.s^{0.6}}$ and $n=0.6$.
Equations~\eqref{eq:moduli} and \eqref{eq:shear-conversion} give
\begin{align}
 G&=\SI{416.6667}{Pa},&\lambda_L&=\SI{277.7778}{Pa},&K_b&=\SI{555.5556}{Pa},\\
 \sigma_y&=\SI{51.96152}{Pa},&&K_e=\SI{24.08225}{Pa.s^{0.6}}.
\end{align}
For an illustrative trial $q_{\rm tr}=\SI{80}{Pa}$ and
$\Delta t=\SI{1.4e-3}{s}$, the bracket endpoints give
$f(\sigma_y)\simeq-\SI{28.0385}{Pa}$ and
$f(q_{\rm tr})\simeq\SI{2.25493}{Pa}$.
The scalar solution is
\begin{align}
 q&\simeq\SI{\HBReturnedStress}{Pa},&
 \Delta\varepsilon_p&\simeq\num{\HBPlasticIncrement},\\
 \dot\varepsilon_p&\simeq\SI{1.139453}{s^{-1}},&
 \alpha&\simeq0.97507446.
\end{align}
The unrounded bracketed calculation satisfies the stated residual tolerance.
Equation~\eqref{eq:plastic-loss} gives
$\Delta D_P/V_0=\SI{\HBDissipation}{J.m^{-3}}$, whereas the isolated return's
elastic-energy decrease is \SI{\HBElasticDrop}{J.m^{-3}}.
Their difference is approximately \SI{0.0015905}{J.m^{-3}}, or 1.28\% of the
recorded plastic-loss estimate for this local example. These numbers do not
represent a complete simulated slump step. The document generator independently
bisects the same scalar equation to make this calculation reproducible.

\paragraph{Implementation mapping.}
\path{core/src/mpm/constitutive.rs}: \nolinkurl{Material::paste},
\nolinkurl{Material::paste_from_shear_rheology},
\nolinkurl{solve_herschel_bulkley_return}, \nolinkurl{symmetric_eigen3} and
\nolinkurl{update_paste}. \path{core/src/mpm/solver.rs} multiplies the returned
loss density by reference point volume and accumulates it only for accepted steps.

\subsection{Weakly compressible liquid: EOS, viscosity, and volume update}
For the liquid wire material \nolinkurl{water}, yield stress is zero and $n=1$;
the consistency field acts as dynamic viscosity $\eta$. The pair $(E,\nu_P)$
is merely a way to specify artificial $K_b=E/[3(1-2\nu_P)]$, not a claim that
liquid has a static Young's modulus or shear elasticity. Define
\begin{equation}
 \rho=\frac{\rho_0}{J},\qquad
 p=K_b\left(\frac1J-1\right)=K_b\left(\frac\rho{\rho_0}-1\right),\qquad
 \mathbf D=\tfrac12(\mathbf C+\mathbf C^T).
\end{equation}
The actual discrete update and stress are
\begin{align}
 J^{n+1}&=J^n[1+\Delta t\operatorname{tr}\mathbf C],\label{eq:liquid-j}\\
 \boldsymbol\tau^{n+1}&=-p^{n+1}J^{n+1}\mathbf I
                  +2\eta J^{n+1}\operatorname{dev}\mathbf D.
 \label{eq:liquid-stress}
\end{align}
The tensor-history slot is reset to identity for liquid; $J$ supplies volume
history. There is no pressure projection, cavitation clamp or pore-pressure
mixture coupling. Expanded states $J>1$ can produce negative EOS pressure.
States with $J^{n+1}\leq10^{-6}$ or nonfinite values are rejected.

Equation~\eqref{eq:liquid-j} is a first-order approximation to
$\dot J=J\operatorname{tr}\nabla\mathbf v$, \emph{not}
$J^n\det(\mathbf I+\Delta t\mathbf C)$. For $\mathbf C=-\mathbf I\,\si{s^{-1}}$,
$J^n=1$ and $\Delta t=\SI{1e-3}{s}$, these expressions give $0.997000000$
and $0.997002999$, respectively. The discrepancy is order $\Delta t^2$ per
step; neither expression is the exact finite-time exponential for constant
$\mathbf C$. No determinant replacement is silently assumed in this derivation.

The liquid elastic-energy density is
\begin{equation}
 \Psi_l(J)=K_b(J-1-\ln J),\qquad
 J\frac{d\Psi_l}{dJ}=K_b(J-1)=-pJ.
\end{equation}
This derivative agrees with the isotropic Kirchhoff stress. Viscous stress is
computed, but its dissipated energy is not accumulated in the current global
energy diagnostic. The continuum viscous power density would be
$2\eta\operatorname{dev}\mathbf D:\operatorname{dev}\mathbf D$ per current volume;
its consistent discrete integral remains a missing ledger term.

\subsection{Hydrostatic initialization and artificial-compressibility scales}
Let $d=H_0-z$ denote depth in the placed current configuration. For liquid,
$dp/dd=\rho g$ and $p(0)=0$ give
\begin{equation}
 \frac{dp}{dd}=\rho_0g\left(1+\frac p{K_b}\right)
 \quad\Longrightarrow\quad
 p(d)=K_b\left(e^{\rho_0gd/K_b}-1\right),\qquad
 J(d)=e^{-\rho_0gd/K_b}. \label{eq:liquid-hydrostatic}
\end{equation}
This is an exact \emph{continuum} hydrostatic solution for the specified EOS.
A placed subcell of current volume $V$ is assigned $V_0=V/J$ and
$m=\rho_0V_0$; assigning $m=\rho_0V$ instead would discard the compressible
density profile. Accurate initial pressure values do not establish discrete
force balance or time-evolved hydrostatic accuracy.

Paste initialization instead uses a confined-compression reference, with
$M=\lambda_L+2G$,
\begin{equation}
 s=-\frac{\rho_0gd}{M},\quad
 \mathbf F_e=\operatorname{diag}(1,1,e^s),\quad
 \boldsymbol\tau=\operatorname{diag}(\lambda_Ls,\lambda_Ls,Ms).
 \label{eq:paste-hydrostatic}
\end{equation}
It matches the \emph{small-strain} vertical reference $-\tau_{zz}=\rho_0gd$.
At finite strain, density is $\rho_0/e^s$ and Cauchy stress is
$\boldsymbol\tau/e^s$, so this is not an exact finite-strain hydrostatic solution.
It is also a prestressed initializer, not an equilibrium solve subject to the
paste yield criterion. For the soft slump values and $H_0=\SI{20}{mm}$,
$\rho_0gH_0/M\simeq0.177$, not an asymptotically small strain.

The acoustic speeds used by the timestep heuristic are
\begin{equation}
 c_{\rm paste}=\sqrt{\frac{K_b+4G/3}{\rho_0}},\qquad
 c_{\rm liquid}=\sqrt{\frac{K_b}{\rho_0}},\qquad
 \nu_{\rm liquid}=\frac\eta{\rho_0}.
\end{equation}
Paste has zero explicit diffusivity in this heuristic because its viscous return
is local/implicit; global kinematics and stress forces still have an acoustic
step restriction. Increasing $K_b$ suppresses artificial compression while
increasing acoustic cost. Useful scale checks are
\begin{equation}
 \epsilon_c=\frac{\rho_0gH_0}{K_b},\qquad
 \mathrm{Ma}=\frac Uc,
\end{equation}
with small $\epsilon_c$ and Mach number desirable for a near-incompressible
approximation, but neither a substitute for boundary and pressure verification.

For $E=\SI{1e5}{Pa}$, $\nu_P=0.2$, $\rho_0=\SI{1000}{kg.m^{-3}}$,
$\eta=\SI{1e-3}{Pa.s}$ and $H_0=\SI{10}{mm}$,
\begin{align}
 K_b&=\SI{\WaterBulkModulus}{Pa},&c&=\SI{\WaterSoundSpeed}{m.s^{-1}},\\
 p(H_0)&=\SI{\WaterBottomPressure}{Pa},&J(H_0)&=\num{\WaterBottomRatio}.
\end{align}
The incompressible comparison is $\rho_0gH_0=\SI{98.1}{Pa}$; the artificial
compression is about 0.176\%. With $h=\SI{2}{mm}$, CFL $=0.3$ and negligible
speed, the acoustic limit is \SI{\WaterAcousticStep}{s}; the viscous limit
$h^2/(4\nu)$ is \SI{1}{s}. A water-like bulk modulus of \SI{2.2e9}{Pa} at
the same reference density would instead give $c\simeq\SI{1483}{m.s^{-1}}$
and an acoustic limit near \SI{4.05e-7}{s}, about 199 times smaller. This cost
tradeoff changes the modeled material; it is not a free numerical acceleration
of identical water physics.

\paragraph{Implementation mapping and unresolved accuracy.}
\path{core/src/mpm/constitutive.rs}: \nolinkurl{Material::liquid},
\nolinkurl{update_liquid}, \nolinkurl{liquid_pressure}, \nolinkurl{wave_speed} and
\nolinkurl{diffusivity}. \path{core/src/mpm/fixtures.rs}:
\nolinkurl{FixtureSpec::hydrostatic_state}, \nolinkurl{hydrostatic_reference} and
\nolinkurl{place_particles}. \path{core/src/mpm/solver.rs}:
\nolinkurl{Simulation::elastic_energy}. Existing audits report unresolved
bottom-boundary hydrostatic error and timestep-insensitive per-particle liquid
pressure noise. Initialization tests, correct EOS algebra and a converged local
return do not resolve those spatial/discrete issues.
